\begin{problem}[第34届物理决赛][薄膜干涉与光压]{87}
    用薄膜制备技术在某均质硅基片上沉积一层均匀等厚氮化镓薄膜，制备出一个硅基氮化镓样品，如图 I 所示。
    \begin{subproblem}
        \begin{subsubproblem}
            当用波长范围为 $450\sim1200\,\mathrm{nm}$ 的光垂直均匀照射该样品氮化镓表面，观察到其反射光谱仅有两种波长的光获得最大相干加强，其中之一波长为 $600\,\mathrm{nm}$；氮化镓的折射率 $n$ 与入射光在真空中波长 $\lambda$（单位 $\mathrm{nm}$）之间的关系（色散关系）为
            $n^{2} = 2.26^{2} + \frac{330.1^{2}}{\lambda^{2} - 265.7^{2}}$。
            硅的折射率随波长在 $3.49\sim5.49$ 范围内变化。只考虑氮化镓表面和氮化镓-硅基片界面的反射，求氮化镓薄膜的厚度 $d$ 和另一种获得最大相干加强光的波长。
        \end{subsubproblem}
        \begin{subsubproblem}
            在该样品硅基片的另一面左、右对称的两个半面上分别均匀沉积一光谱选择性材料涂层，如图 II 所示；对某种特定波长的光，左半面涂层全吸收，右半面全反射。用两根长均为 $a$ 的轻细线竖直悬挂该样品，样品长为 $a$、宽为 $b$，可绕过其中心的光滑竖直固定轴 $OO'$ 转动，也可上下移动，如图 III 所示。开始时，样品静止，用上述特定波长的强激光持续均匀垂直照射该样品涂层表面。此后保持激光方向始终不变，样品绕 $OO'$ 轴转动直至稳定。涂层表面始终被激光完全照射。不计激光对样品侧面的照射。设硅基片厚度为 $d'$、密度为 $\rho'$，氮化镓薄膜的厚度为 $d$、密度为 $\rho$，涂层质量可忽略，真空的介电常量 $\varepsilon_{0}$，重力加速度大小为 $g$。若样品稳定后相对于光照前原位形的转角为 $\alpha$，求所用强激光的电场强度有效值 $E$。
        \end{subsubproblem}
        \begin{subsubproblem}
            取 $E = 5.00\times10^{4}\,\mathrm{V\cdot m^{-1}}$, $d' = 3.00\times10^{-4}\,\mathrm{m}$, $\rho' = 2.33\times10^{3}\,\mathrm{kg\cdot m^{-3}}$, $\rho = 6.10\times10^{3}\,\mathrm{kg\cdot m^{-3}}$, $\varepsilon_{0} = 8.85\times10^{-12}\,\mathrm{F\cdot m^{-1}}$, $g = 9.80\,\mathrm{m\cdot s^{-2}}$，求 $\alpha$ 的值。
        \end{subsubproblem}
    \end{subproblem}
\end{problem}